\begin{prob}
    True or False. Suppose $G = (V, E, \omega)$ is a weighted graph for which
    all edges are positive, except for those edges of a node $s$ which may or
    may not be negative. If Dijkstra's algorithm is run on $G$ with $s$ as
    the source, the correct shortest paths will be found. Assume that the graph does not have any negative loops. Justify your answer.

    \begin{soln}
            True.
            The reason that Dijkstra's algorithm can fail in the presence of
            negative edges is that it cannot rule out that a negative edge that
            it has not seen will cause a path which is currently longer to
            eventually become shorter. But if all of the negative edges are
            attached to the source, the rest of the edges are still positive;
            and so there is no way in which adding edges to a path can decrease
            its length.
    \end{soln}
\end{prob}
